\documentclass[12pt,letterpaper]{article}
\usepackage[margin=0.8in]{geometry}
\usepackage{amsmath,amssymb}
\pagestyle{empty}
\setlength{\parindent}{0pt}
\setlength{\parskip}{0.7em}
\begin{document}

1. An unstable particle has an average lifetime of $10^{-7}$ seconds in its rest frame. A beam of these particles is produced with speed $(4/5)c$. How far on average do the particles travel before decaying? Use $c=3.00\times10^8\,\mathrm{m/s}$.

A. $14.4\,\mathrm{m}$ \qquad
B. $20\,\mathrm{m}$ \qquad
C. $30\,\mathrm{m}$ \qquad
D. $40\,\mathrm{m}$ \qquad
E. $50\,\mathrm{m}$

\vfill
\newpage

2. An accelerator experiment produces a new unstable particle $X$ at rest. $X$ quickly decays to a tau (rest mass $1777\,\mathrm{MeV}/c^2$) and a tau antineutrino (which we will take to be massless). Neutrinos are very difficult to observe, so its energy cannot be measured. The tau is measured to have total energy $4000\,\mathrm{MeV}$. What was the rest mass of the mystery $X$ particle?

\vfill
\newpage

3. A spaceship travels at speed $v=(\sqrt{3}/2)c$ from Earth towards a star ten light-years away. The star is stationary relative to Earth. When the ship is halfway to its destination in the Earth frame, the star goes supernova.

(a) In the ship frame, how far from Earth is the ship when the supernova occurs? The positions are to be evaluated simultaneously in the ship frame.

\vfill

(b) How much time passes in the ship frame between leaving Earth and seeing the light from the supernova?

\vfill
\newpage

4. Two bullets collide head-on. The first bullet is going with speed $0.1c$ and has rest mass $5\,\mathrm{g}$. The second bullet is going with speed $0.2c$ and has rest mass $3\,\mathrm{g}$. Assume that after the collision, the two bullets form one stuck-together mass.

What are the rest mass and velocity of the resulting object? State your positive direction.

\vfill
\newpage

5. A star emits light with wavelength $500\,\mathrm{nm}$ in its own rest frame. Its speed relative to Earth is $v$. Give your answers in terms of $v$ and $c$.

(a) If the star is moving directly away from Earth, what wavelength does an observer on Earth detect?

\vfill

(b) If the star is moving perpendicular to the line of sight instead, what is the observed wavelength? The line of sight is defined in the Earth frame at emission.

\vfill
\newpage

6. A particle of rest mass $10m$ moves at $+0.60c$ along the laboratory $x$-axis. It decays into two particles, each of rest mass $3m$. In the parent's rest frame, daughter A moves along $+x$ and daughter B along $-x$. No other particles are produced.

Express energies in multiples of $mc^2$, momenta in multiples of $mc$, and velocities in multiples of $c$.

(a) Find the total energy and signed $x$-momentum of each daughter in the parent rest frame.

\vfill

(b) Find the total energy and signed $x$-momentum of each daughter in the laboratory frame.

\vfill

(c) Find each daughter's signed laboratory velocity.

\vfill
\newpage

7. Treat the electrons as free and initially stationary. You may use $h/(m_ec)=0.00243\,\mathrm{nm}$.

(a) X-rays of wavelength $0.100\,\mathrm{nm}$ scatter directly backwards off electrons in a thin sheet of metal. What is the wavelength of the scattered photons?

\vfill

(b) An optical photon of wavelength $500\,\mathrm{nm}$ scatters off an electron at an angle of $45^\circ$. How big is the wavelength shift as a fraction of the original wavelength?

\vfill
\newpage

8. Two proposed experiments each begin with two particles of rest mass $m$ and total laboratory energy $6mc^2$.

In experiment A, the particles move towards one another with equal and opposite momenta.

In experiment B, one particle is initially stationary.

In each experiment, the particles combine into a single product, with no radiation or other products.

Which experiment produces the larger product rest mass? Calculate the product rest mass in each case and briefly explain the difference.

\vfill
\end{document}
